\section{Artificial Intelligence for Atmospheric Science}

Artificial Intelligence provides new approaches for atmospheric
prediction, analysis and optimization.



\section{Machine Learning Applications}

ACF integrates:

\begin{itemize}
\item Bias correction
\item Statistical forecasting
\item Pattern recognition
\item Anomaly detection
\end{itemize}



\section{Deep Learning Models}


\subsection{Convolutional Neural Networks}

CNN architectures are used for:

\begin{itemize}
\item Satellite image analysis
\item Radar interpretation
\item Spatial feature extraction
\end{itemize}



\subsection{Recurrent Neural Networks}

RNN and LSTM models analyze:

\begin{itemize}
\item Time series
\item Climate signals
\item Atmospheric cycles
\end{itemize}



\subsection{Transformer Weather Models}

Modern AI weather systems use attention mechanisms:

\begin{equation}
Attention(Q,K,V)
=
softmax
\left(
\frac{QK^T}{\sqrt{d_k}}
\right)V
\end{equation}



\section{Physics-Informed Artificial Intelligence}

ACF combines physical equations with neural networks.

The objective function becomes:

\begin{equation}
Loss =
Loss_{data}
+
\lambda Loss_{physics}
\end{equation}


This enables hybrid Physics-AI atmospheric modelling.
